\chapter{The Simplex Method}

The graphical method solves linear programmes in two variables, but real
planning problems involve many more. The simplex method, devised by George Dantzig in 1947, solves LPs of
any size: it hops from vertex to vertex of the feasible polytope, improving
the objective at every step, and every hop is nothing more than a round of
elementary row operations on a tableau. This chapter drills the standard
maximum case (all constraints of type $\le$ with non-negative right-hand
sides) until choosing the pivot, carrying out the row operations, and
reading off the solution are automatic.

\section{Core Rules}

\begin{rulebox}{Standard maximum form and slack variables}
A \emph{standard maximum problem} reads
\[
  \max z=c^{\top}x \quad\text{s.t.}\quad Ax\le b,\; x\ge 0,
  \qquad\text{with } b\ge 0 .
\]
For each constraint $a_i^{\top}x\le b_i$ introduce a slack variable $s_i=b_i-a_i^{\top}x\ge 0$, turning the system into
equalities:
\[
  Ax+Is=b,\qquad x,s\ge 0 .
\]
A slack measures the unused capacity of its constraint: $s_i=0$ exactly
when constraint $i$ binds. If some $b_i<0$, first multiply that row by
$-1$ (which flips the inequality); mixed $\ge$ or $=$ constraints need the
two-phase or Big-M method and are not treated here.
\end{rulebox}

\begin{rulebox}{The initial tableau}
With $m$ constraints and $n$ decision variables the tableau has one row per
constraint, one column per variable (decision and slack) and a right-hand
side (RHS) column. The objective is rewritten as $-c^{\top}x+z=0$ and placed
in the bottom row, so it starts with the \emph{negatives} of the objective
coefficients. In block form:
\begin{center}
\begin{tabular}{c|cc|c}
Basis & $x$ & $s$ & RHS\\
\hline
$s$ & $A$ & $I$ & $b$\\
\hline
$z$ & $-c^{\top}$ & $0$ & $0$\\
\end{tabular}
\end{center}
The initial basis consists of the slacks; the corresponding
basic feasible solution is $x=0$, $s=b$, $z=0$. We write $R_1,\dots,R_m$
for the constraint rows (top to bottom) and $R_z$ for the objective row.
\end{rulebox}

\begin{rulebox}{Reading a basic solution from a tableau}
At every stage of the algorithm:
\begin{itemize}[leftmargin=1.6em,itemsep=0.3ex]
  \item the \textbf{basic variables} are the row labels; each equals the
        RHS entry of its row (its column is a unit vector);
  \item all other (\textbf{non-basic}) variables equal $0$;
  \item the current objective value $z$ is the RHS entry of the $z$-row.
\end{itemize}
A slack that is basic and positive means its constraint holds with room to
spare; a slack equal to zero means its constraint binds.
\end{rulebox}

\begin{rulebox}{Pivot column: the entering variable}
Inspect the $z$-row. The pivot column is the column with
the \textbf{most negative} entry; its variable \emph{enters} the basis,
because each unit of it raises $z$ fastest among all candidates. Ties may
be broken arbitrarily (e.g.\ leftmost column). If the $z$-row contains
\textbf{no negative entry}, stop: the current tableau is optimal.
\end{rulebox}

\begin{rulebox}{Pivot row: the minimum-ratio test}
Only rows with a \textbf{strictly positive} entry in the pivot column take
part; for each such row compute
\[
  \text{ratio}=\frac{\text{RHS}}{\text{entry in pivot column}} .
\]
The row with the \textbf{smallest ratio} is the pivot row; its label is the \emph{leaving} variable, and the entry
where pivot row and pivot column cross is the \textbf{pivot element}. Rows
with entry $\le 0$ are skipped: increasing the entering variable never
drives their basic variable negative. If \emph{no} entry in the pivot
column is positive, the LP is unbounded above.
\end{rulebox}

\begin{rulebox}{The pivot operation (recipe)}
Let $a$ be the pivot element in pivot row $R_p$.
\begin{enumerate}[leftmargin=1.8em,itemsep=0.3ex]
  \item \textbf{Normalise:} $R_p\to R_p/a$, so the pivot element becomes
        $1$.
  \item \textbf{Clear the column:} for every other row, including $R_z$,
        subtract the appropriate multiple of the \emph{new} pivot row:
        $R_i\to R_i-(\text{entry of }R_i\text{ in pivot column})\cdot R_p$.
        Rows with entry $0$ in the pivot column are unchanged.
  \item \textbf{Relabel:} the entering variable replaces the leaving
        variable as the label of row $p$.
\end{enumerate}
Afterwards the pivot column is a unit vector and the new basic solution can
be read off as before.
\end{rulebox}

\begin{rulebox}{A worked pivot}
Take $\max z=5x_1+4x_2$ s.t.\ $6x_1+4x_2\le 24$, $x_1+2x_2\le 6$,
$x_1,x_2\ge 0$. Initial tableau:
\begin{center}
\begin{tabular}{c|cccc|c}
Basis & $x_1$ & $x_2$ & $s_1$ & $s_2$ & RHS\\
\hline
$s_1$ & $6$ & $4$ & $1$ & $0$ & $24$\\
$s_2$ & $1$ & $2$ & $0$ & $1$ & $6$\\
\hline
$z$ & $-5$ & $-4$ & $0$ & $0$ & $0$\\
\end{tabular}
\end{center}
Pivot column $x_1$ ($-5$ is most negative). Ratios: $24/6=4$ and $6/1=6$;
the minimum is $4$, so $s_1$ leaves and the pivot element is $6$.
Row operations: $R_1\to R_1/6$; then $R_2\to R_2-1\cdot R_1$ and
$R_z\to R_z+5R_1$ (both with the \emph{new} $R_1$):
\begin{center}
\begin{tabular}{c|cccc|c}
Basis & $x_1$ & $x_2$ & $s_1$ & $s_2$ & RHS\\
\hline
$x_1$ & $1$ & $\tfrac{2}{3}$ & $\tfrac{1}{6}$ & $0$ & $4$\\
$s_2$ & $0$ & $\tfrac{4}{3}$ & $-\tfrac{1}{6}$ & $1$ & $2$\\
\hline
$z$ & $0$ & $-\tfrac{2}{3}$ & $\tfrac{5}{6}$ & $0$ & $20$\\
\end{tabular}
\end{center}
New basic solution: $x_1=4$, $s_2=2$, $x_2=s_1=0$, $z=20$. The $z$-row
still contains $-\tfrac{2}{3}$, so another iteration would follow.
\end{rulebox}

\begin{rulebox}{Stopping rule and optimum}
Iterate (pivot column, ratio test, pivot) until the $z$-row has no
negative entry. Then the tableau is optimal: the basic
variables take their RHS values, non-basic variables are zero, and the
bottom-right entry is $z^{*}$. Every pivot moves to an adjacent vertex of
the feasible polytope without ever decreasing $z$.
\end{rulebox}

\section{Worked Example}

\begin{worked}{Run the simplex method on the running example, from the initial tableau to the optimum}
\wstep{Step 1. The initial tableau.}
With the standard form of Chapter~3 and the objective written as
$z-40x_{1}-30x_{2}=0$:
\begin{center}\footnotesize
\begin{tabular}{c|rrrr|r}
  \hline
  basis & $x_{1}$ & $x_{2}$ & $s_{1}$ & $s_{2}$ & RHS\\
  \hline
  $s_{1}$ & $2$ & $1$ & $1$ & $0$ & $100$\\
  $s_{2}$ & $1$ & $1$ & $0$ & $1$ & $80$\\
  \hline
  $z$ & $-40$ & $-30$ & $0$ & $0$ & $0$\\
  \hline
\end{tabular}
\end{center}
The basis is $\{s_{1},s_{2}\}$, the basic solution is $(0,0,100,80)$, and
$z=0$: produce nothing, all capacity idle. That is the corner $(0,0)$.
\wtrap{With the objective written as $z-\mathbf c^{\!\top}\mathbf x=0$, the
$z$-row holds \emph{negated} objective coefficients. Signs in the $z$-row
are the single largest source of confusion in this chapter; fix the
convention --- \emph{negative entries mean improvement is available} --- and
never change it mid-problem.}

\wstep{Step 2. First iteration: choose the entering variable.}
The most negative $z$-row entry is $-40$, under $x_{1}$. Bringing $x_{1}$
into the basis raises $z$ by $40$ per unit, faster than $x_{2}$'s $30$. So
$x_{1}$ \textbf{enters}.

\wstep{Step 3. Minimum-ratio test for the leaving variable.}
Divide each RHS by the \emph{positive} entries in the pivot column:
\[
  \text{row }s_{1}:\ \frac{100}{2}=50 ;
  \qquad
  \text{row }s_{2}:\ \frac{80}{1}=80 .
\]
The smallest ratio, $50$, is in row $s_{1}$: that row's resource runs out
first, so $s_{1}$ \textbf{leaves}. The pivot element is the $2$.
\wtrap{Use only \emph{strictly positive} entries in the ratio test. A zero
or negative entry imposes no limit --- increasing the entering variable
never drives that basic variable negative --- and including it produces a
nonsensical ratio and an infeasible next tableau. If \emph{no} entry is
positive, the problem is unbounded (Chapter~6).}

\wstep{Step 4. Pivot: make the pivot element $1$ and clear its column.}
Divide the pivot row by $2$, then subtract multiples of it from the others:
\begin{center}\footnotesize
\begin{tabular}{c|rrrr|r}
  \hline
  basis & $x_{1}$ & $x_{2}$ & $s_{1}$ & $s_{2}$ & RHS\\
  \hline
  $x_{1}$ & $1$ & $0.5$ & $0.5$ & $0$ & $50$\\
  $s_{2}$ & $0$ & $0.5$ & $-0.5$ & $1$ & $30$\\
  \hline
  $z$ & $0$ & $-10$ & $20$ & $0$ & $2000$\\
  \hline
\end{tabular}
\end{center}
Basic solution $(50,0,0,30)$, $z=2000$ --- the corner $(50,0)$. The method
has moved along an edge of the feasible region.
\wtrap{Clear the pivot column in \emph{every} other row, the $z$-row
included. Forgetting the $z$-row leaves a stale objective and the algorithm
loses track of where it is. After each pivot, the pivot column should be a
unit vector.}

\wstep{Step 5. Second iteration.}
$-10$ under $x_{2}$ is still negative, so $x_{2}$ enters. Ratios:
\[
  \text{row }x_{1}:\ \frac{50}{0.5}=100 ;
  \qquad
  \text{row }s_{2}:\ \frac{30}{0.5}=60 .
\]
Minimum $60$, so $s_{2}$ leaves; the pivot element is the $0.5$ in the
$s_{2}$ row.

\wstep{Step 6. Pivot again.}
Multiply that row by $2$, then clear the $x_{2}$ column:
\begin{center}\footnotesize
\begin{tabular}{c|rrrr|r}
  \hline
  basis & $x_{1}$ & $x_{2}$ & $s_{1}$ & $s_{2}$ & RHS\\
  \hline
  $x_{1}$ & $1$ & $0$ & $1$ & $-1$ & $20$\\
  $x_{2}$ & $0$ & $1$ & $-1$ & $2$ & $60$\\
  \hline
  $z$ & $0$ & $0$ & $10$ & $20$ & $2600$\\
  \hline
\end{tabular}
\end{center}

\wstep{Step 7. Stop, and read everything off.}
No negative entry remains in the $z$-row, so no variable can improve $z$:
\textbf{optimal}.
\[
  \boxed{x_{1}^{*}=20,\quad x_{2}^{*}=60,\quad s_{1}^{*}=s_{2}^{*}=0,\quad z^{*}=2600 .}
\]
This matches the graphical answer of Chapter~2 exactly. \checkmark\ The path
taken was $(0,0)\to(50,0)\to(20,60)$: two pivots, two edges, three corners.

\wstep{Step 8. The $z$-row entries under the slacks are not decoration.}
The $10$ under $s_{1}$ and $20$ under $s_{2}$ are the \textbf{shadow
prices}: one more machine hour is worth $10$ euros, one more labour hour
$20$. Chapter~5 reads them systematically and Chapter~7 proves they are the
dual solution.

\wstep{Step 9. The algorithm in five lines.}
\begin{center}\footnotesize
\begin{tabular}{ll}
  \hline
  1 & any negative $z$-row entry? If not, stop --- optimal.\\
  2 & entering: the most negative one (its column is the pivot column).\\
  3 & leaving: minimum ratio RHS $/$ positive column entry.\\
  4 & pivot: normalise the pivot row, clear the column everywhere else.\\
  5 & go to 1.\\
  \hline
\end{tabular}
\end{center}
Each pass moves to an adjacent corner with a value at least as good, and
since there are finitely many corners the method terminates.
\end{worked}

\section{Practice Problems}

All variables are non-negative in every problem. For each complete simplex
run show the initial tableau, every intermediate tableau, and state the row
operations of each pivot.

\begin{problems}
  \item Rewrite $\max z=2x_1+3x_2$ s.t.\ $x_1+x_2\le 4$, $2x_1+x_2\le 6$
        in equation form with slack variables, set up the initial tableau,
        and read off the initial basic feasible solution.
  \item Solve the LP of Problem~1 with the simplex method.
  \item Solve $\max z=3x_1+2x_2$ s.t.\ $x_1+x_2\le 4$, $x_1+3x_2\le 6$.
  \item Solve $\max z=3x_1+5x_2$ s.t.\ $x_1\le 4$, $2x_2\le 12$,
        $3x_1+2x_2\le 18$.
  \item The initial tableau of $\max z=4x_1+3x_2$ s.t.\ $2x_1+x_2\le 8$,
        $x_1+3x_2\le 9$ is
        \begin{center}
        \begin{tabular}{c|cccc|c}
        Basis & $x_1$ & $x_2$ & $s_1$ & $s_2$ & RHS\\
        \hline
        $s_1$ & $2$ & $1$ & $1$ & $0$ & $8$\\
        $s_2$ & $1$ & $3$ & $0$ & $1$ & $9$\\
        \hline
        $z$ & $-4$ & $-3$ & $0$ & $0$ & $0$\\
        \end{tabular}
        \end{center}
        Determine the entering variable, all ratios of the minimum-ratio
        test, the leaving variable, and the pivot element. Do \emph{not}
        carry out the pivot.
  \item Perform exactly one simplex iteration on $\max z=3x_1+x_2$ s.t.\
        $x_1+x_2\le 4$, $2x_1+x_2\le 6$. Is the resulting tableau optimal?
  \item A simplex run for $\max z=2x_1+5x_2$ s.t.\ $x_1+2x_2\le 10$,
        $x_2\le 4$ ended with the tableau
        \begin{center}
        \begin{tabular}{c|cccc|c}
        Basis & $x_1$ & $x_2$ & $s_1$ & $s_2$ & RHS\\
        \hline
        $x_1$ & $1$ & $0$ & $1$ & $-2$ & $2$\\
        $x_2$ & $0$ & $1$ & $0$ & $1$ & $4$\\
        \hline
        $z$ & $0$ & $0$ & $2$ & $1$ & $24$\\
        \end{tabular}
        \end{center}
        Verify that it is optimal, read off the optimal solution and
        objective value, and decide which constraints bind.
  \item Solve $\max z=2x_1+x_2+x_3$ s.t.\ $x_1+x_2+x_3\le 6$,
        $2x_1+x_2\le 8$.
  \item Solve $\max z=4x_1+5x_2$ s.t.\ $2x_1+x_2\le 8$, $x_1+2x_2\le 7$.
  \item The initial tableau of $\max z=2x_1+3x_2$ s.t.\ $3x_1-x_2\le 6$,
        $x_1+2x_2\le 4$ is
        \begin{center}
        \begin{tabular}{c|cccc|c}
        Basis & $x_1$ & $x_2$ & $s_1$ & $s_2$ & RHS\\
        \hline
        $s_1$ & $3$ & $-1$ & $1$ & $0$ & $6$\\
        $s_2$ & $1$ & $2$ & $0$ & $1$ & $4$\\
        \hline
        $z$ & $-2$ & $-3$ & $0$ & $0$ & $0$\\
        \end{tabular}
        \end{center}
        Determine the pivot column and the pivot row, and explain why one
        of the rows does not take part in the minimum-ratio test.
  \item Solve $\max z=3x_1+2x_2$ s.t.\ $x_1+x_2\le 4$, $x_1+2x_2\le 6$,
        $x_1\le 3$.
  \item Solve $\max z=5x_1+4x_2+3x_3$ s.t.\ $2x_1+3x_2+x_3\le 5$,
        $4x_1+x_2+2x_3\le 11$, $3x_1+4x_2+2x_3\le 8$.
  \item Solve $\max z=x_1+2x_2$ s.t.\ $x_1+x_2\le 5$, $x_1\le 3$.
  \item For $\max z=6x_1+5x_2$ s.t.\ $x_1+x_2\le 5$, $3x_1+2x_2\le 12$
        one simplex iteration has already been carried out, producing
        \begin{center}
        \begin{tabular}{c|cccc|c}
        Basis & $x_1$ & $x_2$ & $s_1$ & $s_2$ & RHS\\
        \hline
        $s_1$ & $0$ & $\tfrac{1}{3}$ & $1$ & $-\tfrac{1}{3}$ & $1$\\
        $x_1$ & $1$ & $\tfrac{2}{3}$ & $0$ & $\tfrac{1}{3}$ & $4$\\
        \hline
        $z$ & $0$ & $-1$ & $0$ & $2$ & $24$\\
        \end{tabular}
        \end{center}
        (a)~Read off the current basic solution. (b)~Is the tableau
        optimal? (c)~If not, iterate to the optimum.
\end{problems}

\section{Solutions}

\begin{solutions}
  \item With slacks $s_1,s_2\ge 0$ the constraints become
        $x_1+x_2+s_1=4$ and $2x_1+x_2+s_2=6$; the objective row is
        $-2x_1-3x_2+z=0$. Initial tableau:
        \begin{center}
        \begin{tabular}{c|cccc|c}
        Basis & $x_1$ & $x_2$ & $s_1$ & $s_2$ & RHS\\
        \hline
        $s_1$ & $1$ & $1$ & $1$ & $0$ & $4$\\
        $s_2$ & $2$ & $1$ & $0$ & $1$ & $6$\\
        \hline
        $z$ & $-2$ & $-3$ & $0$ & $0$ & $0$\\
        \end{tabular}
        \end{center}
        Initial basic feasible solution: $x_1=x_2=0$ (non-basic),
        $s_1=4$, $s_2=6$, $z=0$.

  \item Start from the tableau of Solution~1. Pivot column: $x_2$
        ($-3$ is the most negative $z$-row entry). Ratios: $R_1$:
        $4/1=4$; $R_2$: $6/1=6$. Minimum $4$, so $s_1$ leaves; pivot
        element $1$. Row operations: $R_1$ unchanged (division by $1$),
        relabel $s_1\to x_2$; $R_2\to R_2-R_1$; $R_z\to R_z+3R_1$.
        \begin{center}
        \begin{tabular}{c|cccc|c}
        Basis & $x_1$ & $x_2$ & $s_1$ & $s_2$ & RHS\\
        \hline
        $x_2$ & $1$ & $1$ & $1$ & $0$ & $4$\\
        $s_2$ & $1$ & $0$ & $-1$ & $1$ & $2$\\
        \hline
        $z$ & $1$ & $0$ & $3$ & $0$ & $12$\\
        \end{tabular}
        \end{center}
        No negative entry in the $z$-row: optimal. $x_1^{*}=0$,
        $x_2^{*}=4$, $z^{*}=12$; $s_2=2$, so constraint 1 binds and
        constraint 2 has $2$ units of unused capacity.

  \item Initial tableau:
        \begin{center}
        \begin{tabular}{c|cccc|c}
        Basis & $x_1$ & $x_2$ & $s_1$ & $s_2$ & RHS\\
        \hline
        $s_1$ & $1$ & $1$ & $1$ & $0$ & $4$\\
        $s_2$ & $1$ & $3$ & $0$ & $1$ & $6$\\
        \hline
        $z$ & $-3$ & $-2$ & $0$ & $0$ & $0$\\
        \end{tabular}
        \end{center}
        Pivot column $x_1$ ($-3$). Ratios: $4/1=4$, $6/1=6$; pivot row
        $R_1$, $s_1$ leaves, pivot element $1$. Row operations: relabel
        $s_1\to x_1$; $R_2\to R_2-R_1$; $R_z\to R_z+3R_1$.
        \begin{center}
        \begin{tabular}{c|cccc|c}
        Basis & $x_1$ & $x_2$ & $s_1$ & $s_2$ & RHS\\
        \hline
        $x_1$ & $1$ & $1$ & $1$ & $0$ & $4$\\
        $s_2$ & $0$ & $2$ & $-1$ & $1$ & $2$\\
        \hline
        $z$ & $0$ & $1$ & $3$ & $0$ & $12$\\
        \end{tabular}
        \end{center}
        All $z$-row entries non-negative: optimal. $x_1^{*}=4$,
        $x_2^{*}=0$, $z^{*}=12$; constraint 1 binds, constraint 2 keeps
        slack $s_2=2$.

  \item Initial tableau:
        \begin{center}
        \begin{tabular}{c|ccccc|c}
        Basis & $x_1$ & $x_2$ & $s_1$ & $s_2$ & $s_3$ & RHS\\
        \hline
        $s_1$ & $1$ & $0$ & $1$ & $0$ & $0$ & $4$\\
        $s_2$ & $0$ & $2$ & $0$ & $1$ & $0$ & $12$\\
        $s_3$ & $3$ & $2$ & $0$ & $0$ & $1$ & $18$\\
        \hline
        $z$ & $-3$ & $-5$ & $0$ & $0$ & $0$ & $0$\\
        \end{tabular}
        \end{center}
        \emph{Iteration 1.} Pivot column $x_2$ ($-5$). Ratios: $R_1$ is
        skipped (entry $0$); $R_2$: $12/2=6$; $R_3$: $18/2=9$. Pivot row
        $R_2$, $s_2$ leaves, pivot element $2$. Row operations:
        $R_2\to R_2/2$; then $R_3\to R_3-2R_2$, $R_z\to R_z+5R_2$;
        $R_1$ unchanged.
        \begin{center}
        \begin{tabular}{c|ccccc|c}
        Basis & $x_1$ & $x_2$ & $s_1$ & $s_2$ & $s_3$ & RHS\\
        \hline
        $s_1$ & $1$ & $0$ & $1$ & $0$ & $0$ & $4$\\
        $x_2$ & $0$ & $1$ & $0$ & $\tfrac{1}{2}$ & $0$ & $6$\\
        $s_3$ & $3$ & $0$ & $0$ & $-1$ & $1$ & $6$\\
        \hline
        $z$ & $-3$ & $0$ & $0$ & $\tfrac{5}{2}$ & $0$ & $30$\\
        \end{tabular}
        \end{center}
        \emph{Iteration 2.} Pivot column $x_1$ ($-3$). Ratios: $R_1$:
        $4/1=4$; $R_2$ skipped (entry $0$); $R_3$: $6/3=2$. Pivot row
        $R_3$, $s_3$ leaves, pivot element $3$. Row operations:
        $R_3\to R_3/3$; then $R_1\to R_1-R_3$, $R_z\to R_z+3R_3$.
        \begin{center}
        \begin{tabular}{c|ccccc|c}
        Basis & $x_1$ & $x_2$ & $s_1$ & $s_2$ & $s_3$ & RHS\\
        \hline
        $s_1$ & $0$ & $0$ & $1$ & $\tfrac{1}{3}$ & $-\tfrac{1}{3}$ & $2$\\
        $x_2$ & $0$ & $1$ & $0$ & $\tfrac{1}{2}$ & $0$ & $6$\\
        $x_1$ & $1$ & $0$ & $0$ & $-\tfrac{1}{3}$ & $\tfrac{1}{3}$ & $2$\\
        \hline
        $z$ & $0$ & $0$ & $0$ & $\tfrac{3}{2}$ & $1$ & $36$\\
        \end{tabular}
        \end{center}
        Optimal: $x_1^{*}=2$, $x_2^{*}=6$, $z^{*}=36$; $s_1=2$, so
        constraint 1 is slack while constraints 2 and 3 bind.

  \item Entering variable: $x_1$, because $-4$ is the most negative
        $z$-row entry. Ratios: $R_1$: $8/2=4$; $R_2$: $9/1=9$. The
        minimum is $4$, so $s_1$ is the leaving variable and the pivot
        element is the $2$ at the crossing of row $s_1$ and column
        $x_1$.

  \item Initial tableau:
        \begin{center}
        \begin{tabular}{c|cccc|c}
        Basis & $x_1$ & $x_2$ & $s_1$ & $s_2$ & RHS\\
        \hline
        $s_1$ & $1$ & $1$ & $1$ & $0$ & $4$\\
        $s_2$ & $2$ & $1$ & $0$ & $1$ & $6$\\
        \hline
        $z$ & $-3$ & $-1$ & $0$ & $0$ & $0$\\
        \end{tabular}
        \end{center}
        Pivot column $x_1$ ($-3$). Ratios: $4/1=4$, $6/2=3$; pivot row
        $R_2$, $s_2$ leaves, pivot element $2$. Row operations:
        $R_2\to R_2/2$; then $R_1\to R_1-R_2$, $R_z\to R_z+3R_2$.
        \begin{center}
        \begin{tabular}{c|cccc|c}
        Basis & $x_1$ & $x_2$ & $s_1$ & $s_2$ & RHS\\
        \hline
        $s_1$ & $0$ & $\tfrac{1}{2}$ & $1$ & $-\tfrac{1}{2}$ & $1$\\
        $x_1$ & $1$ & $\tfrac{1}{2}$ & $0$ & $\tfrac{1}{2}$ & $3$\\
        \hline
        $z$ & $0$ & $\tfrac{1}{2}$ & $0$ & $\tfrac{3}{2}$ & $9$\\
        \end{tabular}
        \end{center}
        All $z$-row entries are non-negative, so this single iteration
        already reaches the optimum: $x_1^{*}=3$, $x_2^{*}=0$,
        $z^{*}=9$, with $s_1=1$.

  \item The $z$-row $(0,0,2,1\mid 24)$ has no negative entry, so the
        tableau is optimal. Basic variables (row labels): $x_1=2$,
        $x_2=4$; non-basic: $s_1=s_2=0$; $z^{*}=24$. Since both slacks
        are zero, both constraints bind. Check: $2+2\cdot 4=10$ and
        $x_2=4$; $z=2\cdot 2+5\cdot 4=24$.

  \item Initial tableau:
        \begin{center}
        \begin{tabular}{c|ccccc|c}
        Basis & $x_1$ & $x_2$ & $x_3$ & $s_1$ & $s_2$ & RHS\\
        \hline
        $s_1$ & $1$ & $1$ & $1$ & $1$ & $0$ & $6$\\
        $s_2$ & $2$ & $1$ & $0$ & $0$ & $1$ & $8$\\
        \hline
        $z$ & $-2$ & $-1$ & $-1$ & $0$ & $0$ & $0$\\
        \end{tabular}
        \end{center}
        \emph{Iteration 1.} Pivot column $x_1$ ($-2$). Ratios: $6/1=6$,
        $8/2=4$; pivot row $R_2$, $s_2$ leaves, pivot element $2$. Row
        operations: $R_2\to R_2/2$; then $R_1\to R_1-R_2$,
        $R_z\to R_z+2R_2$.
        \begin{center}
        \begin{tabular}{c|ccccc|c}
        Basis & $x_1$ & $x_2$ & $x_3$ & $s_1$ & $s_2$ & RHS\\
        \hline
        $s_1$ & $0$ & $\tfrac{1}{2}$ & $1$ & $1$ & $-\tfrac{1}{2}$ & $2$\\
        $x_1$ & $1$ & $\tfrac{1}{2}$ & $0$ & $0$ & $\tfrac{1}{2}$ & $4$\\
        \hline
        $z$ & $0$ & $0$ & $-1$ & $0$ & $1$ & $8$\\
        \end{tabular}
        \end{center}
        \emph{Iteration 2.} Pivot column $x_3$ ($-1$). Only $R_1$ has a
        positive entry ($1$); ratio $2/1=2$; $s_1$ leaves, pivot element
        $1$. Row operations: $R_1$ unchanged, relabel $s_1\to x_3$;
        $R_2$ unchanged (entry $0$); $R_z\to R_z+R_1$.
        \begin{center}
        \begin{tabular}{c|ccccc|c}
        Basis & $x_1$ & $x_2$ & $x_3$ & $s_1$ & $s_2$ & RHS\\
        \hline
        $x_3$ & $0$ & $\tfrac{1}{2}$ & $1$ & $1$ & $-\tfrac{1}{2}$ & $2$\\
        $x_1$ & $1$ & $\tfrac{1}{2}$ & $0$ & $0$ & $\tfrac{1}{2}$ & $4$\\
        \hline
        $z$ & $0$ & $\tfrac{1}{2}$ & $0$ & $1$ & $\tfrac{1}{2}$ & $10$\\
        \end{tabular}
        \end{center}
        Optimal: $x_1^{*}=4$, $x_2^{*}=0$, $x_3^{*}=2$, $z^{*}=10$;
        both slacks are non-basic, so both constraints bind.

  \item Initial tableau:
        \begin{center}
        \begin{tabular}{c|cccc|c}
        Basis & $x_1$ & $x_2$ & $s_1$ & $s_2$ & RHS\\
        \hline
        $s_1$ & $2$ & $1$ & $1$ & $0$ & $8$\\
        $s_2$ & $1$ & $2$ & $0$ & $1$ & $7$\\
        \hline
        $z$ & $-4$ & $-5$ & $0$ & $0$ & $0$\\
        \end{tabular}
        \end{center}
        \emph{Iteration 1.} Pivot column $x_2$ ($-5$). Ratios: $8/1=8$,
        $7/2=\tfrac{7}{2}$; pivot row $R_2$, $s_2$ leaves, pivot element
        $2$. Row operations: $R_2\to R_2/2$; then $R_1\to R_1-R_2$,
        $R_z\to R_z+5R_2$.
        \begin{center}
        \begin{tabular}{c|cccc|c}
        Basis & $x_1$ & $x_2$ & $s_1$ & $s_2$ & RHS\\
        \hline
        $s_1$ & $\tfrac{3}{2}$ & $0$ & $1$ & $-\tfrac{1}{2}$ & $\tfrac{9}{2}$\\
        $x_2$ & $\tfrac{1}{2}$ & $1$ & $0$ & $\tfrac{1}{2}$ & $\tfrac{7}{2}$\\
        \hline
        $z$ & $-\tfrac{3}{2}$ & $0$ & $0$ & $\tfrac{5}{2}$ & $\tfrac{35}{2}$\\
        \end{tabular}
        \end{center}
        \emph{Iteration 2.} Pivot column $x_1$ ($-\tfrac{3}{2}$).
        Ratios: $R_1$: $\tfrac{9}{2}\big/\tfrac{3}{2}=3$; $R_2$:
        $\tfrac{7}{2}\big/\tfrac{1}{2}=7$. Pivot row $R_1$, $s_1$
        leaves, pivot element $\tfrac{3}{2}$. Row operations:
        $R_1\to \tfrac{2}{3}R_1$; then $R_2\to R_2-\tfrac{1}{2}R_1$,
        $R_z\to R_z+\tfrac{3}{2}R_1$.
        \begin{center}
        \begin{tabular}{c|cccc|c}
        Basis & $x_1$ & $x_2$ & $s_1$ & $s_2$ & RHS\\
        \hline
        $x_1$ & $1$ & $0$ & $\tfrac{2}{3}$ & $-\tfrac{1}{3}$ & $3$\\
        $x_2$ & $0$ & $1$ & $-\tfrac{1}{3}$ & $\tfrac{2}{3}$ & $2$\\
        \hline
        $z$ & $0$ & $0$ & $1$ & $2$ & $22$\\
        \end{tabular}
        \end{center}
        Optimal: $x_1^{*}=3$, $x_2^{*}=2$, $z^{*}=22$; both constraints
        bind ($6+2=8$, $3+4=7$).

  \item Pivot column: $x_2$, since $-3$ is the most negative $z$-row
        entry. In the ratio test row $R_1$ is excluded because its entry
        in the pivot column, $-1$, is not strictly positive: increasing
        $x_2$ only \emph{loosens} constraint 1, so $s_1$ can never be
        driven negative. The only ratio is $R_2$: $4/2=2$, so the pivot
        row is $R_2$, $s_2$ leaves, and the pivot element is $2$.

  \item Initial tableau:
        \begin{center}
        \begin{tabular}{c|ccccc|c}
        Basis & $x_1$ & $x_2$ & $s_1$ & $s_2$ & $s_3$ & RHS\\
        \hline
        $s_1$ & $1$ & $1$ & $1$ & $0$ & $0$ & $4$\\
        $s_2$ & $1$ & $2$ & $0$ & $1$ & $0$ & $6$\\
        $s_3$ & $1$ & $0$ & $0$ & $0$ & $1$ & $3$\\
        \hline
        $z$ & $-3$ & $-2$ & $0$ & $0$ & $0$ & $0$\\
        \end{tabular}
        \end{center}
        \emph{Iteration 1.} Pivot column $x_1$ ($-3$). Ratios:
        $4/1=4$, $6/1=6$, $3/1=3$; pivot row $R_3$, $s_3$ leaves, pivot
        element $1$. Row operations: relabel $s_3\to x_1$;
        $R_1\to R_1-R_3$; $R_2\to R_2-R_3$; $R_z\to R_z+3R_3$.
        \begin{center}
        \begin{tabular}{c|ccccc|c}
        Basis & $x_1$ & $x_2$ & $s_1$ & $s_2$ & $s_3$ & RHS\\
        \hline
        $s_1$ & $0$ & $1$ & $1$ & $0$ & $-1$ & $1$\\
        $s_2$ & $0$ & $2$ & $0$ & $1$ & $-1$ & $3$\\
        $x_1$ & $1$ & $0$ & $0$ & $0$ & $1$ & $3$\\
        \hline
        $z$ & $0$ & $-2$ & $0$ & $0$ & $3$ & $9$\\
        \end{tabular}
        \end{center}
        \emph{Iteration 2.} Pivot column $x_2$ ($-2$). Ratios: $R_1$:
        $1/1=1$; $R_2$: $3/2=\tfrac{3}{2}$; $R_3$ skipped (entry $0$).
        Pivot row $R_1$, $s_1$ leaves, pivot element $1$. Row
        operations: relabel $s_1\to x_2$; $R_2\to R_2-2R_1$;
        $R_z\to R_z+2R_1$; $R_3$ unchanged.
        \begin{center}
        \begin{tabular}{c|ccccc|c}
        Basis & $x_1$ & $x_2$ & $s_1$ & $s_2$ & $s_3$ & RHS\\
        \hline
        $x_2$ & $0$ & $1$ & $1$ & $0$ & $-1$ & $1$\\
        $s_2$ & $0$ & $0$ & $-2$ & $1$ & $1$ & $1$\\
        $x_1$ & $1$ & $0$ & $0$ & $0$ & $1$ & $3$\\
        \hline
        $z$ & $0$ & $0$ & $2$ & $0$ & $1$ & $11$\\
        \end{tabular}
        \end{center}
        Optimal: $x_1^{*}=3$, $x_2^{*}=1$, $z^{*}=11$; $s_2=1$, so
        constraints 1 and 3 bind while constraint 2 is slack.

  \item Initial tableau:
        \begin{center}
        \begin{tabular}{c|cccccc|c}
        Basis & $x_1$ & $x_2$ & $x_3$ & $s_1$ & $s_2$ & $s_3$ & RHS\\
        \hline
        $s_1$ & $2$ & $3$ & $1$ & $1$ & $0$ & $0$ & $5$\\
        $s_2$ & $4$ & $1$ & $2$ & $0$ & $1$ & $0$ & $11$\\
        $s_3$ & $3$ & $4$ & $2$ & $0$ & $0$ & $1$ & $8$\\
        \hline
        $z$ & $-5$ & $-4$ & $-3$ & $0$ & $0$ & $0$ & $0$\\
        \end{tabular}
        \end{center}
        \emph{Iteration 1.} Pivot column $x_1$ ($-5$). Ratios:
        $5/2=\tfrac{5}{2}$, $11/4=\tfrac{11}{4}$, $8/3=\tfrac{8}{3}$;
        the smallest is $\tfrac{5}{2}$, so $s_1$ leaves and the pivot
        element is $2$. Row operations: $R_1\to R_1/2$; then
        $R_2\to R_2-4R_1$, $R_3\to R_3-3R_1$, $R_z\to R_z+5R_1$.
        \begin{center}
        \begin{tabular}{c|cccccc|c}
        Basis & $x_1$ & $x_2$ & $x_3$ & $s_1$ & $s_2$ & $s_3$ & RHS\\
        \hline
        $x_1$ & $1$ & $\tfrac{3}{2}$ & $\tfrac{1}{2}$ & $\tfrac{1}{2}$ & $0$ & $0$ & $\tfrac{5}{2}$\\
        $s_2$ & $0$ & $-5$ & $0$ & $-2$ & $1$ & $0$ & $1$\\
        $s_3$ & $0$ & $-\tfrac{1}{2}$ & $\tfrac{1}{2}$ & $-\tfrac{3}{2}$ & $0$ & $1$ & $\tfrac{1}{2}$\\
        \hline
        $z$ & $0$ & $\tfrac{7}{2}$ & $-\tfrac{1}{2}$ & $\tfrac{5}{2}$ & $0$ & $0$ & $\tfrac{25}{2}$\\
        \end{tabular}
        \end{center}
        \emph{Iteration 2.} Pivot column $x_3$ ($-\tfrac{1}{2}$).
        Ratios: $R_1$: $\tfrac{5}{2}\big/\tfrac{1}{2}=5$; $R_2$ skipped
        (entry $0$); $R_3$: $\tfrac{1}{2}\big/\tfrac{1}{2}=1$. Pivot row
        $R_3$, $s_3$ leaves, pivot element $\tfrac{1}{2}$. Row
        operations: $R_3\to 2R_3$; then $R_1\to R_1-\tfrac{1}{2}R_3$,
        $R_z\to R_z+\tfrac{1}{2}R_3$; $R_2$ unchanged.
        \begin{center}
        \begin{tabular}{c|cccccc|c}
        Basis & $x_1$ & $x_2$ & $x_3$ & $s_1$ & $s_2$ & $s_3$ & RHS\\
        \hline
        $x_1$ & $1$ & $2$ & $0$ & $2$ & $0$ & $-1$ & $2$\\
        $s_2$ & $0$ & $-5$ & $0$ & $-2$ & $1$ & $0$ & $1$\\
        $x_3$ & $0$ & $-1$ & $1$ & $-3$ & $0$ & $2$ & $1$\\
        \hline
        $z$ & $0$ & $3$ & $0$ & $1$ & $0$ & $1$ & $13$\\
        \end{tabular}
        \end{center}
        Optimal: $x_1^{*}=2$, $x_2^{*}=0$, $x_3^{*}=1$, $z^{*}=13$;
        $s_2=1$, so constraints 1 and 3 bind and constraint 2 has one
        unit of slack. Check: $2\cdot 2+1=5$, $4\cdot 2+2=10\le 11$,
        $3\cdot 2+2=8$; $z=10+3=13$.

  \item Initial tableau:
        \begin{center}
        \begin{tabular}{c|cccc|c}
        Basis & $x_1$ & $x_2$ & $s_1$ & $s_2$ & RHS\\
        \hline
        $s_1$ & $1$ & $1$ & $1$ & $0$ & $5$\\
        $s_2$ & $1$ & $0$ & $0$ & $1$ & $3$\\
        \hline
        $z$ & $-1$ & $-2$ & $0$ & $0$ & $0$\\
        \end{tabular}
        \end{center}
        Pivot column $x_2$ ($-2$). Ratios: $R_1$: $5/1=5$; $R_2$ is
        skipped (entry $0$: the bound $x_1\le 3$ does not restrict
        $x_2$). Pivot row $R_1$, $s_1$ leaves, pivot element $1$. Row
        operations: relabel $s_1\to x_2$; $R_2$ unchanged;
        $R_z\to R_z+2R_1$.
        \begin{center}
        \begin{tabular}{c|cccc|c}
        Basis & $x_1$ & $x_2$ & $s_1$ & $s_2$ & RHS\\
        \hline
        $x_2$ & $1$ & $1$ & $1$ & $0$ & $5$\\
        $s_2$ & $1$ & $0$ & $0$ & $1$ & $3$\\
        \hline
        $z$ & $1$ & $0$ & $2$ & $0$ & $10$\\
        \end{tabular}
        \end{center}
        The $x_1$ entry of the $z$-row became $-1+2\cdot 1=1\ge 0$, so
        the tableau is optimal after one pivot: $x_1^{*}=0$,
        $x_2^{*}=5$, $z^{*}=10$; $s_2=3$.

  \item (a)~Basic variables: $s_1=1$ and $x_1=4$; non-basic:
        $x_2=s_2=0$; current objective $z=24$. (b)~Not optimal: the
        $x_2$ column still carries $-1$ in the $z$-row. (c)~Pivot column
        $x_2$. Ratios: $R_1$: $1\big/\tfrac{1}{3}=3$; $R_2$:
        $4\big/\tfrac{2}{3}=6$. Pivot row $R_1$, $s_1$ leaves, pivot
        element $\tfrac{1}{3}$. Row operations: $R_1\to 3R_1$; then
        $R_2\to R_2-\tfrac{2}{3}R_1$, $R_z\to R_z+R_1$.
        \begin{center}
        \begin{tabular}{c|cccc|c}
        Basis & $x_1$ & $x_2$ & $s_1$ & $s_2$ & RHS\\
        \hline
        $x_2$ & $0$ & $1$ & $3$ & $-1$ & $3$\\
        $x_1$ & $1$ & $0$ & $-2$ & $1$ & $2$\\
        \hline
        $z$ & $0$ & $0$ & $3$ & $1$ & $27$\\
        \end{tabular}
        \end{center}
        Optimal: $x_1^{*}=2$, $x_2^{*}=3$, $z^{*}=27$; both slacks are
        zero, so both constraints bind ($2+3=5$, $6+6=12$).
\end{solutions}

\section{Additional Problems}

As before: all variables are non-negative, and complete runs must show
every tableau and every row operation.

\begin{problems}
  \item Solve $\max z=2x_1+x_2$ s.t.\ $x_1+x_2\le 3$, $x_1+2x_2\le 4$.
  \item Solve $\max z=4x_1+6x_2$ s.t.\ $x_1+x_2\le 5$, $x_1+3x_2\le 9$.
  \item The initial tableau of $\max z=5x_1+6x_2$ s.t.\ $x_1+4x_2\le 8$,
        $2x_1+3x_2\le 12$ is
        \begin{center}
        \begin{tabular}{c|cccc|c}
        Basis & $x_1$ & $x_2$ & $s_1$ & $s_2$ & RHS\\
        \hline
        $s_1$ & $1$ & $4$ & $1$ & $0$ & $8$\\
        $s_2$ & $2$ & $3$ & $0$ & $1$ & $12$\\
        \hline
        $z$ & $-5$ & $-6$ & $0$ & $0$ & $0$\\
        \end{tabular}
        \end{center}
        Determine the entering variable, all ratios, the leaving
        variable, and the pivot element. Do \emph{not} pivot.
  \item For $\max z=4x_1+3x_2$ s.t.\ $2x_1+x_2\le 6$, $x_1+3x_2\le 9$,
        set up the initial tableau, perform the first pivot, read off the
        new basic solution, and decide whether the new tableau is
        optimal.
  \item Solve $\max z=3x_1+2x_2+x_3$ s.t.\ $x_1+x_2+x_3\le 5$,
        $2x_1+x_2\le 6$.
  \item Solve $\max z=4x_1+3x_2$ s.t.\ $2x_1+x_2\le 10$, $x_1+x_2\le 6$,
        $x_1+2x_2\le 10$.
\end{problems}

\section{Additional Solutions}

\begin{solutions}
  \item Initial tableau:
        \begin{center}
        \begin{tabular}{c|cccc|c}
        Basis & $x_1$ & $x_2$ & $s_1$ & $s_2$ & RHS\\
        \hline
        $s_1$ & $1$ & $1$ & $1$ & $0$ & $3$\\
        $s_2$ & $1$ & $2$ & $0$ & $1$ & $4$\\
        \hline
        $z$ & $-2$ & $-1$ & $0$ & $0$ & $0$\\
        \end{tabular}
        \end{center}
        Pivot column $x_1$ ($-2$). Ratios: $3/1=3$, $4/1=4$; pivot row
        $R_1$, $s_1$ leaves, pivot element $1$. Row operations: relabel
        $s_1\to x_1$; $R_2\to R_2-R_1$; $R_z\to R_z+2R_1$.
        \begin{center}
        \begin{tabular}{c|cccc|c}
        Basis & $x_1$ & $x_2$ & $s_1$ & $s_2$ & RHS\\
        \hline
        $x_1$ & $1$ & $1$ & $1$ & $0$ & $3$\\
        $s_2$ & $0$ & $1$ & $-1$ & $1$ & $1$\\
        \hline
        $z$ & $0$ & $1$ & $2$ & $0$ & $6$\\
        \end{tabular}
        \end{center}
        Optimal: $x_1^{*}=3$, $x_2^{*}=0$, $z^{*}=6$; $s_2=1$, so only
        constraint 1 binds.

  \item Initial tableau:
        \begin{center}
        \begin{tabular}{c|cccc|c}
        Basis & $x_1$ & $x_2$ & $s_1$ & $s_2$ & RHS\\
        \hline
        $s_1$ & $1$ & $1$ & $1$ & $0$ & $5$\\
        $s_2$ & $1$ & $3$ & $0$ & $1$ & $9$\\
        \hline
        $z$ & $-4$ & $-6$ & $0$ & $0$ & $0$\\
        \end{tabular}
        \end{center}
        \emph{Iteration 1.} Pivot column $x_2$ ($-6$). Ratios: $5/1=5$,
        $9/3=3$; pivot row $R_2$, $s_2$ leaves, pivot element $3$. Row
        operations: $R_2\to R_2/3$; then $R_1\to R_1-R_2$,
        $R_z\to R_z+6R_2$.
        \begin{center}
        \begin{tabular}{c|cccc|c}
        Basis & $x_1$ & $x_2$ & $s_1$ & $s_2$ & RHS\\
        \hline
        $s_1$ & $\tfrac{2}{3}$ & $0$ & $1$ & $-\tfrac{1}{3}$ & $2$\\
        $x_2$ & $\tfrac{1}{3}$ & $1$ & $0$ & $\tfrac{1}{3}$ & $3$\\
        \hline
        $z$ & $-2$ & $0$ & $0$ & $2$ & $18$\\
        \end{tabular}
        \end{center}
        \emph{Iteration 2.} Pivot column $x_1$ ($-2$). Ratios: $R_1$:
        $2\big/\tfrac{2}{3}=3$; $R_2$: $3\big/\tfrac{1}{3}=9$. Pivot row
        $R_1$, $s_1$ leaves, pivot element $\tfrac{2}{3}$. Row
        operations: $R_1\to \tfrac{3}{2}R_1$; then
        $R_2\to R_2-\tfrac{1}{3}R_1$, $R_z\to R_z+2R_1$.
        \begin{center}
        \begin{tabular}{c|cccc|c}
        Basis & $x_1$ & $x_2$ & $s_1$ & $s_2$ & RHS\\
        \hline
        $x_1$ & $1$ & $0$ & $\tfrac{3}{2}$ & $-\tfrac{1}{2}$ & $3$\\
        $x_2$ & $0$ & $1$ & $-\tfrac{1}{2}$ & $\tfrac{1}{2}$ & $2$\\
        \hline
        $z$ & $0$ & $0$ & $3$ & $1$ & $24$\\
        \end{tabular}
        \end{center}
        Optimal: $x_1^{*}=3$, $x_2^{*}=2$, $z^{*}=24$; both constraints
        bind ($3+2=5$, $3+6=9$).

  \item Entering variable: $x_2$ ($-6$ is most negative). Ratios:
        $R_1$: $8/4=2$; $R_2$: $12/3=4$. The minimum is $2$, so $s_1$
        leaves and the pivot element is the $4$ in row $s_1$, column
        $x_2$.

  \item Initial tableau:
        \begin{center}
        \begin{tabular}{c|cccc|c}
        Basis & $x_1$ & $x_2$ & $s_1$ & $s_2$ & RHS\\
        \hline
        $s_1$ & $2$ & $1$ & $1$ & $0$ & $6$\\
        $s_2$ & $1$ & $3$ & $0$ & $1$ & $9$\\
        \hline
        $z$ & $-4$ & $-3$ & $0$ & $0$ & $0$\\
        \end{tabular}
        \end{center}
        Pivot column $x_1$ ($-4$). Ratios: $6/2=3$, $9/1=9$; pivot row
        $R_1$, $s_1$ leaves, pivot element $2$. Row operations:
        $R_1\to R_1/2$; then $R_2\to R_2-R_1$, $R_z\to R_z+4R_1$.
        \begin{center}
        \begin{tabular}{c|cccc|c}
        Basis & $x_1$ & $x_2$ & $s_1$ & $s_2$ & RHS\\
        \hline
        $x_1$ & $1$ & $\tfrac{1}{2}$ & $\tfrac{1}{2}$ & $0$ & $3$\\
        $s_2$ & $0$ & $\tfrac{5}{2}$ & $-\tfrac{1}{2}$ & $1$ & $6$\\
        \hline
        $z$ & $0$ & $-1$ & $2$ & $0$ & $12$\\
        \end{tabular}
        \end{center}
        New basic solution: $x_1=3$, $s_2=6$, $x_2=s_1=0$, $z=12$. The
        tableau is \emph{not} optimal: the $x_2$ column still has $-1$
        in the $z$-row, so a further pivot (with $x_2$ entering) would
        follow.

  \item Initial tableau:
        \begin{center}
        \begin{tabular}{c|ccccc|c}
        Basis & $x_1$ & $x_2$ & $x_3$ & $s_1$ & $s_2$ & RHS\\
        \hline
        $s_1$ & $1$ & $1$ & $1$ & $1$ & $0$ & $5$\\
        $s_2$ & $2$ & $1$ & $0$ & $0$ & $1$ & $6$\\
        \hline
        $z$ & $-3$ & $-2$ & $-1$ & $0$ & $0$ & $0$\\
        \end{tabular}
        \end{center}
        \emph{Iteration 1.} Pivot column $x_1$ ($-3$). Ratios: $5/1=5$,
        $6/2=3$; pivot row $R_2$, $s_2$ leaves, pivot element $2$. Row
        operations: $R_2\to R_2/2$; then $R_1\to R_1-R_2$,
        $R_z\to R_z+3R_2$.
        \begin{center}
        \begin{tabular}{c|ccccc|c}
        Basis & $x_1$ & $x_2$ & $x_3$ & $s_1$ & $s_2$ & RHS\\
        \hline
        $s_1$ & $0$ & $\tfrac{1}{2}$ & $1$ & $1$ & $-\tfrac{1}{2}$ & $2$\\
        $x_1$ & $1$ & $\tfrac{1}{2}$ & $0$ & $0$ & $\tfrac{1}{2}$ & $3$\\
        \hline
        $z$ & $0$ & $-\tfrac{1}{2}$ & $-1$ & $0$ & $\tfrac{3}{2}$ & $9$\\
        \end{tabular}
        \end{center}
        \emph{Iteration 2.} Pivot column $x_3$ ($-1$ beats
        $-\tfrac{1}{2}$). Only $R_1$ has a positive entry ($1$); ratio
        $2/1=2$; $s_1$ leaves, pivot element $1$. Row operations:
        relabel $s_1\to x_3$; $R_2$ unchanged (entry $0$);
        $R_z\to R_z+R_1$.
        \begin{center}
        \begin{tabular}{c|ccccc|c}
        Basis & $x_1$ & $x_2$ & $x_3$ & $s_1$ & $s_2$ & RHS\\
        \hline
        $x_3$ & $0$ & $\tfrac{1}{2}$ & $1$ & $1$ & $-\tfrac{1}{2}$ & $2$\\
        $x_1$ & $1$ & $\tfrac{1}{2}$ & $0$ & $0$ & $\tfrac{1}{2}$ & $3$\\
        \hline
        $z$ & $0$ & $0$ & $0$ & $1$ & $1$ & $11$\\
        \end{tabular}
        \end{center}
        Optimal: $x_1^{*}=3$, $x_2^{*}=0$, $x_3^{*}=2$, $z^{*}=11$;
        both constraints bind ($3+0+2=5$, $2\cdot 3+0=6$).

  \item Initial tableau:
        \begin{center}
        \begin{tabular}{c|ccccc|c}
        Basis & $x_1$ & $x_2$ & $s_1$ & $s_2$ & $s_3$ & RHS\\
        \hline
        $s_1$ & $2$ & $1$ & $1$ & $0$ & $0$ & $10$\\
        $s_2$ & $1$ & $1$ & $0$ & $1$ & $0$ & $6$\\
        $s_3$ & $1$ & $2$ & $0$ & $0$ & $1$ & $10$\\
        \hline
        $z$ & $-4$ & $-3$ & $0$ & $0$ & $0$ & $0$\\
        \end{tabular}
        \end{center}
        \emph{Iteration 1.} Pivot column $x_1$ ($-4$). Ratios:
        $10/2=5$, $6/1=6$, $10/1=10$; pivot row $R_1$, $s_1$ leaves,
        pivot element $2$. Row operations: $R_1\to R_1/2$; then
        $R_2\to R_2-R_1$, $R_3\to R_3-R_1$, $R_z\to R_z+4R_1$.
        \begin{center}
        \begin{tabular}{c|ccccc|c}
        Basis & $x_1$ & $x_2$ & $s_1$ & $s_2$ & $s_3$ & RHS\\
        \hline
        $x_1$ & $1$ & $\tfrac{1}{2}$ & $\tfrac{1}{2}$ & $0$ & $0$ & $5$\\
        $s_2$ & $0$ & $\tfrac{1}{2}$ & $-\tfrac{1}{2}$ & $1$ & $0$ & $1$\\
        $s_3$ & $0$ & $\tfrac{3}{2}$ & $-\tfrac{1}{2}$ & $0$ & $1$ & $5$\\
        \hline
        $z$ & $0$ & $-1$ & $2$ & $0$ & $0$ & $20$\\
        \end{tabular}
        \end{center}
        \emph{Iteration 2.} Pivot column $x_2$ ($-1$). Ratios: $R_1$:
        $5\big/\tfrac{1}{2}=10$; $R_2$: $1\big/\tfrac{1}{2}=2$; $R_3$:
        $5\big/\tfrac{3}{2}=\tfrac{10}{3}$. Pivot row $R_2$, $s_2$
        leaves, pivot element $\tfrac{1}{2}$. Row operations:
        $R_2\to 2R_2$; then $R_1\to R_1-\tfrac{1}{2}R_2$,
        $R_3\to R_3-\tfrac{3}{2}R_2$, $R_z\to R_z+R_2$.
        \begin{center}
        \begin{tabular}{c|ccccc|c}
        Basis & $x_1$ & $x_2$ & $s_1$ & $s_2$ & $s_3$ & RHS\\
        \hline
        $x_1$ & $1$ & $0$ & $1$ & $-1$ & $0$ & $4$\\
        $x_2$ & $0$ & $1$ & $-1$ & $2$ & $0$ & $2$\\
        $s_3$ & $0$ & $0$ & $1$ & $-3$ & $1$ & $2$\\
        \hline
        $z$ & $0$ & $0$ & $1$ & $2$ & $0$ & $22$\\
        \end{tabular}
        \end{center}
        Optimal: $x_1^{*}=4$, $x_2^{*}=2$, $z^{*}=22$; $s_3=2$, so
        constraints 1 and 2 bind ($8+2=10$, $4+2=6$) while constraint 3
        keeps $2$ units of slack ($4+4=8\le 10$).
\end{solutions}
